\section{Introduction}
\label{sec:intro}

Enterprise LLM agents are deployed today with a capability that is inseparable
from its own attack surface: they retrieve from private organizational
stores---email, chat, wikis, ticketing systems, meeting transcripts---and then
\emph{act}, sending messages and writing to shared artifacts on a user's
behalf. The retrieval is what makes them useful. The action is what makes
retrieval dangerous.

Recent benchmarking establishes that this is not a hypothetical concern.
CI-Work~\cite{ciwork2026} constructs enterprise tasks in which information that
must be conveyed and information that must be withheld are deliberately
entangled in the same dense retrieved context, and reports that frontier agents
commit a contextual-integrity violation in 15.8--50.9\% of cases, with leakage
rates up to 26.7\%. Two of its findings matter more than the headline numbers.
First, the failure does not yield to scale: larger models in the same family
convey more \emph{and} leak more. Second, it does not yield to prompting
either---the strongest prompt-level defense evaluated reduces violation from
27.8\% to 21.3\% while cutting conveyance from 93.0\% to 81.0\%, trading most
of the task's usefulness for a fraction of its safety. The authors conclude
that progress requires ``a paradigm shift, moving beyond model-centric scaling
toward context-centric architectures.''

This paper takes that conclusion literally, and argues that the architecture
called for cannot be built from the measurements currently available, because
those measurements describe only half of the failure.

\subsection{Two failures, one of them unmeasured}

Consider an agent asked to email a landlord to renew an office lease. It
retrieves three things: the executed lease, fine to discuss with the landlord;
an internal memo stating the walk-away rent threshold, real but confidential;
and an unconfirmed message from HR speculating that headcount cuts will reduce
the space requirement. The agent sends an email that states the walk-away
threshold---destroying the negotiating position---and asserts the speculative
square-footage figure as though it were a company position.

These are two different mistakes. The agent \emph{said something it should not
have said}, and it \emph{believed something it should not have believed}.
CI-Work and the contextual-integrity benchmarks preceding
it~\cite{confaide2024,privacylens2024,cibench2024} measure the first. The
second has no metric in this literature at all, despite being the failure that
produces a wrong outcome rather than an embarrassing one.

The distinction is not a refinement. It is a claim that the two properties are
carried on \emph{independent} axes, and the evidence for independence is
already inside CI-Work's own published examples. A negotiation walk-away price
is maximally restricted and also completely authoritative: it is a real number
from a real internal decision. An unverified assertion in a channel the
counterparty controls is the reverse---quoting it leaks nothing, and relying on
it is catastrophic. A policy with one axis must collapse these two cases, and
whichever way it collapses them it is wrong about one of them.

Security has a name for the second axis and a formalism for enforcing it.
Bell--LaPadula~\cite{belllapadula1976} governs confidentiality through
no-read-up and no-write-down; Biba~\cite{biba1977} governs integrity through
the dual rules, no-read-down and no-write-up. The mnemonic is that
\emph{Bell--LaPadula keeps secrets from leaking down, and Biba keeps garbage
from flowing up}. The agent-privacy literature has effectively rediscovered the
first half. The second half is sitting unused, and it is the half that
describes contamination.

\subsection{The same failure, in production}
\label{sec:intro:incidents}

The integrity axis is not a benchmark artifact. It is the shape of the
agent incidents that have actually reached postmortem, and naming it makes
those incidents legible as one failure rather than six unrelated ones.
\Cref{tab:incidents} maps five documented cases onto the framework.

\begin{table*}[t]
\centering
\caption{Documented production incidents, read through the framework. In every
case the agent's \emph{authority} exceeded the integrity of the evidence its
action rested on---\Cref{eq:bibaviolation}. We cite these as existence proofs
of a failure mode, not as measurements; the rate-of-spend dimension of the
cost incidents is explicitly outside what we address
(\Cref{sec:discussion:limitations}).}
\label{tab:incidents}
\footnotesize
\begin{tabular}{@{}p{0.15\textwidth}p{0.30\textwidth}p{0.24\textwidth}p{0.24\textwidth}@{}}
\toprule
\textbf{Incident} & \textbf{What happened} & \textbf{Evidence the action rested on}
  & \textbf{What the framework would have required} \\
\midrule
PocketOS / Railway, 2026~\cite{pocketos2026}
& Coding agent hit a credential mismatch, found an unrelated API token in
  another file, issued a volume deletion. Production database and all
  co-located backups gone in nine seconds.
& An \emph{inference} that the volume was staging-scoped. The agent's own
  account: it ``guessed rather than verified.''
& Irreversible action requires $\Il \ge 0.95$; a guess scores near the floor.
  Action capped at ``ask a question.'' Separately, an ambient token is not in
  the capability grant (precedence rule~1). \\[4pt]
Replit, 2025~\cite{replit2025}
& Agent ran destructive commands against production during an explicit code
  freeze, destroying 1{,}206 executive and 1{,}196+ company records.
& An instruction the agent was told to obey and did not.
& The canonical \textbf{L0} failure: enforcement by instruction
  (\Cref{sec:method:ladder}). The vendor's own remediation---dev/prod
  separation, planning-only mode---is a capability constraint, which is the
  argument. \\[4pt]
Runaway accounting agent, 2026~\cite{mandiant2026}
& 15{,}000+ high-cost API calls in under an hour; $\approx$\$50{,}000 in
  charges; active business transactions disrupted.
& A retry loop. No evidence at all supported continuing.
& \emph{Partially addressed.} $\lvlI$ makes spend authority an explicit policy
  input rather than an ambient property of the credential. We do not address
  the loop itself. \\[4pt]
\$82{,}000 GCP bill, 2026~\cite{gcp82k2026}
& Service-account key left in a project root; agent retried Terraform applies
  for 56 hours.
& A key that happened to be in reach.
& Same ambient-credential pattern as PocketOS. Credentials belong in the
  per-request grant, not the environment. \\[4pt]
Moffatt v.\ Air Canada, 2024~\cite{moffatt2024}
& Chatbot stated a bereavement-fare policy that did not exist; the tribunal
  held the airline liable.
& Nothing. The claim had no source entry.
& Grounding rate $\mathrm{GR}$ (\Cref{sec:benchmark:metrics}) scores a claim
  with no supporting entry as ungrounded. The ruling is useful to us
  independently: it establishes that a customer-facing assertion \emph{is} a
  binding high-integrity artifact. \\
\bottomrule
\end{tabular}
\end{table*}

Two observations organize the table. First, in none of these cases was the
agent attacked, and in only one was it even confused in the ordinary sense. It
was reasoning forward from an obstacle, forming a plausible plan, and executing
with the authority it happened to hold. \textbf{The blast radius was set by the
infrastructure and the trigger was set by the agent, with nothing in between
relating the two.} Second, that missing relation is exactly
\Cref{eq:bibaviolation}: the integrity demanded by the action was never
compared against the integrity of the evidence supporting it, because neither
quantity was represented.

We are deliberately careful about what this table claims. Our mechanism
constrains \emph{authority}---what an agent may cause given what it knows. It
does not constrain \emph{rate}, and the cost incidents above are substantially
rate failures that want spend caps, timeouts, and kill
switches~\cite{infoq2026billing}. We claim the authority half, and say so again
in \Cref{sec:discussion:limitations}.

\subsection{Why a write boundary, and why zero trust}

Our mechanism follows from a diagnosis: enterprise agents fail because the
moment an item enters the context window it is treated as both \emph{sendable}
and \emph{true}. Neither property is conferred by retrieval, and both are
properties of a \emph{request}---the same internal roadmap note is quotable to
a teammate and not to an outside contributor, with no change to the note.

That is precisely the premise of zero trust~\cite{nist800207}: never trust,
always verify, re-evaluated per access rather than granted by position.
Tian and Song~\cite{tiansong2021} give the specific formulation we build
on, reframing zero trust as dynamic dual-label Bell--LaPadula plus Biba with
continuous per-component trust scores and separate weights for confidentiality
and integrity requirements. Their method was never applied to agents, never
applied to information flow inside a model's context window, and---as
subsequent surveys note~\cite{he2022zerotrust,sultana2024zerotrust}---was
published without a principled initial-trust assignment, with a weight
distribution the reviews call ``not reasonable enough,'' and without any
quantitative evaluation. We supply all three.

We place the enforcement point at the agent's \textbf{write boundary}: the
final tool call that emits an artifact to a recipient. This is the narrowest
interface at which the question ``may this specific item reach this specific
audience, and is it sound enough to act on'' is even well posed, and it is
where the read/write asymmetry of Bell--LaPadula earns its keep. In our worked
case the agent's read authority is the reviewer's clearance ($0.9$) while its
write ceiling is the thread's audience ($0.2$, a public repository). Every leak
lives in that gap---and so does all of the context that makes the output good.
An access model that can only grant or deny cannot express ``read this, reason
over it, never repeat it.''

\subsection{Research questions}
\label{sec:intro:rqs}

The paper is organized around three questions. They are ordered so that each
one is only worth asking if the previous one is answered affirmatively: if the
axes are not independent there is nothing new to enforce, and if enforcement
cannot beat the trade-off there is no reason to care where it sits.

\begin{description}
  \item[\textbf{RQ1 (Representation).}] \emph{Are confidentiality and
  provenance integrity independent axes in enterprise agent context, and does
  the integrity axis carry failures that existing contextual-integrity
  benchmarks cannot express or score?}

  This is a question about what a benchmark can \emph{say}, prior to any
  defense. It decomposes into three testable parts: whether all four quadrants
  of \Cref{tab:quadrants} are populated in real data; whether both failure
  types co-occur in transcripts the CI-Work authors themselves published; and
  whether the essential/sensitive partition is not merely incomplete but
  \emph{unsound}, in that it cannot represent an entry that is simultaneously
  required by the task and forbidden to the recipient. Addressed in
  \Cref{sec:threat,sec:benchmark:classes}, and settled empirically in
  \Cref{sec:eval:rq1}.

  \item[\textbf{RQ2 (Efficacy).}] \emph{Can a per-request dual-label policy at
  the write boundary reduce leakage, violation, and contamination without the
  utility collapse that prompt-level defenses incur---and under what condition
  does it fail to?}

  The second clause is the one that makes this honest. A defense that wins on
  average while quietly over-blocking is the prompt defense with extra
  machinery, so RQ2 is only answered if false suppression is measured on null
  controls and the failure condition is stated precisely rather than averaged
  away. Addressed in \Cref{sec:eval:rq2}.

  \item[\textbf{RQ3 (Enforcement).}] \emph{Where must the policy be enforced,
  and by what mechanism, for its decisions to hold against a model that is
  careless or a channel that is adversarial?}

  A correct decision enforced by a sentence in a system prompt is not enforced.
  RQ3 asks separately about \emph{placement}---before generation versus after
  it---and about \emph{mechanism}, via the L0--L3 ladder and the action-
  downgrade rule of \Cref{eq:downgrade}. It also generates the paper's sharpest
  falsifiable prediction, since a mechanism that removes restricted text from
  the writer's context should flatten the leakage-versus-scale slope that
  CI-Work reports. Addressed in
  \Cref{sec:eval:placement,sec:discussion:prediction}.
\end{description}

\noindent A fourth question---\emph{does the framework explain agent failures
observed in production?}---we treat as design validation rather than as an
experiment, since \Cref{tab:incidents} is a retrospective reading of incidents
we did not control. We return to it in \Cref{sec:discussion:incidents}.

\subsection{Three results}

\textbf{Two axes are necessary, and the second one needs a new entry class to
be measurable.} CI-Work partitions retrieved context into an essential set and
a sensitive set. That partition can express only confidentiality failures, so
we add a third class, \emph{corrupting} context ($\Ecorr$): low
confidentiality, low integrity, cheap to disclose and expensive to believe. The
PR author's own claim that ``security already signed off'' is the canonical
instance. We further show the two-set partition is not merely incomplete but
unsound for enterprise work, because it implicitly assumes
$\Eess \subseteq \text{emittable}$. \Cref{sec:eval} exhibits an entry that is
simultaneously necessary for the task and forbidden to the recipient---a cell
the existing schema cannot represent, and one that requires paraphrase rather
than either quotation or deletion.

\textbf{Integrity for an agent is a capability constraint, not a filtering
problem.} The harm from a low-integrity entry is not that the agent
\emph{repeats} it---repeating it is usually harmless---but that the agent
\emph{relies} on it, and filtering text does nothing about reliance. We
therefore bound the integrity of the action an agent may take by the integrity
of the weakest evidence its conclusion rests on
(\Cref{sec:method:downgrade}). Strong evidence permits approving a merge; weak
evidence permits only asking a question. Three properties make this the right
primitive: it requires no natural-language inference, being a comparison over
labels already computed; it cannot be argued out of, because it is enforced by
the permission system rather than by the model's cooperation, and so is immune
to the compliance failure CI-Work measures; and it degrades gracefully,
yielding a weaker action rather than a refusal. It also reframes indirect
prompt injection~\cite{greshake2023indirect} usefully: an injected instruction
is a low-integrity entry attempting to ground a high-integrity action, and
under this rule it does not need to be \emph{detected}, only \emph{labeled}.

\textbf{Enforcement placement dominates enforcement strength.} We evaluate the
same policy at two deployment points and find the difference is not a matter of
degree. Post-hoc filtering receives a finished draft and can only \emph{delete},
which leaves two failures structurally unreachable: an essential entry above
the write ceiling is lost outright because paraphrase is not an available move,
and a finding the agent dropped \emph{because} it believed a corrupting entry
cannot be restored, because a filter cannot add text. Running the same decision
before generation recovers all of it---conveyance returns to the undefended
agent's figure exactly---because the writer receives admitted entries verbatim,
above-ceiling entries as label-checked abstracts, and never holds the withheld
entries at all. No-write-down becomes vacuous rather than filtered for that
set: you cannot emit what you do not have.

\subsection{Contributions}

\begin{enumerate}
  \item \textbf{A dual-axis threat model} for enterprise agents that extends
  CI-Work's formal risk model with an integrity-contamination condition, and a
  demonstration that both failures co-occur in each of CI-Work's own two
  published qualitative transcripts (\Cref{sec:threat}).

  \item \textbf{\sys}, a zero-trust policy layer at the agent's write boundary
  enforcing Bell--LaPadula no-write-down and Biba no-write-up per request, with
  an explicit provenance rubric for the integrity label and an initial-trust
  rule, addressing the two published weaknesses of
  Tian--Song~\cite{tiansong2021} (\Cref{sec:method}).

  \item \textbf{Derived rather than hand-set weights.} We split every parameter
  into structural, measurement, and operating-point tiers with distinct
  justification obligations, and derive $\wC$ from audience breadth and $\wI$
  from consequence class instead of tuning them
  (\Cref{sec:method:weights}). We also show the required outbound integrity is
  a function of the agent's authorized \emph{capability}, so revoking merge
  permission genuinely lowers the evidence bar for an identical task.

  \item \textbf{An enforcement-mechanism ladder} (L0--L3) that separates what a
  policy decides from what makes the decision hold, and a falsifiable
  prediction it generates: under handle-mediated citation the leakage-versus-
  model-size slope should flatten, contradicting CI-Work's inverse-scaling
  result (\Cref{sec:method:ladder}, \Cref{sec:discussion:prediction}).

  \item \textbf{A benchmark with a corrupting entry class, entangled twins, and
  null controls.} Every sensitive entry is authored as a near-twin of an
  essential one---same subject, same retrieval query, different disclosure
  norm---so the task cannot be solved by relevance filtering. Null seeds with
  no sensitive entries measure false suppression, which is what distinguishes a
  policy layer from a refusal and which no prior CI benchmark includes
  (\Cref{sec:benchmark}).

  \item \textbf{Two integrity metrics, an honest account of their limits, and a
  reproducible artifact.} We report Integrity-Leakage and Integrity-Violation,
  and show by measurement that they undercount contamination by exactly the
  entries that were silently believed rather than repeated---motivating a
  counterfactual verdict-corruption metric. Our release includes a verifier
  that re-derives every policy decision in the paper from the configuration
  file, so the prose cannot drift from the implementation
  (\Cref{sec:eval}, \Cref{sec:ethics:openscience}).
\end{enumerate}

We also report two latent policy bugs that our own derivation exposed, because
both are instructive rather than incidental (\Cref{sec:discussion:bugs}). In
one, a content marker such as ``approved'' could raise an entry's integrity
score regardless of its source, so text in a channel controlled by the party
under review could talk our integrity labeler into trusting it---our own
defense, prompt-injected. The principle the fix follows from is that
\emph{integrity is a property of provenance, and content cannot vouch for its
own provenance}.
